\section{Compute Accounting：FLOPs、FLOP/s 与 MFU}

\subsection{FLOPs 与 FLOP/s}

\begin{knowledgebox}{课堂提示：两个同音缩写必须分账}
老师专门提醒 FLOPs 与 FLOP/s 是“terribly confusing acronyms”：前者是做了多少浮点运算，后者是硬件每秒完成多少浮点运算。训练成本预测先算总 FLOPs，性能测量再算 actual FLOP/s；把两者混用，会把“任务规模很大”和“实现效率很低”误判成同一个问题。
\end{knowledgebox}

\begin{longtable}{p{0.18\textwidth}p{0.36\textwidth}p{0.36\textwidth}}
\toprule
\textbf{术语} & \textbf{含义} & \textbf{例子} \\
\midrule
FLOPs & floating-point operations，计算总量 & GPT-3 训练约 \(3.14\times 10^{23}\) FLOPs。 \\
FLOP/s & 每秒浮点运算数，硬件或程序速度 & H100 dense bf16 峰值约 \(989.5\) TFLOP/s。 \\
MFU & actual FLOP/s / promised FLOP/s & 模型实际吃满硬件的程度，0.5 已经很好。 \\
\bottomrule
\end{longtable}

\[
\mathrm{MFU}=\frac{\text{actual FLOP/s}}{\text{promised FLOP/s}}.
\]

\begin{warningbox}{FLOPs 和 FLOP/s 混淆会毁掉估算}
FLOPs 是路程，FLOP/s 是速度。训练时间约等于总 FLOPs 除以有效 FLOP/s。有效速度又取决于 dtype、kernel、通信、shape 和 MFU。
\end{warningbox}

\subsection{矩阵乘法 FLOPs}

前面定义了 FLOP/s 与 MFU，本节用最常见的矩阵乘法建立可复用计数模板。若 \(X\in\mathbb{R}^{B\times D}\)，\(W\in\mathbb{R}^{D\times K}\)，每个输出元素需要沿 $D$ 维做乘加，输出共有 $BK$ 个元素，因此

\[
Y=XW,\qquad
\text{FLOPs}\approx 2BDK.
\]

每个输出元素需要 \(D\) 次乘法和约 \(D\) 次加法。实际程序还要 benchmark 得到运行时间，再算实际 FLOP/s。

\begin{importantbox}{MFU 的正确使用}
MFU 不是道德评分，而是诊断工具。MFU 低可能是 memory-bound、通信同步、kernel 太小、shape 不友好、数据加载慢，也可能是测量口径不一致。下一节 arithmetic intensity 会解释其中最常见的一类原因。
\end{importantbox}

\begin{knowledgebox}{实践经验：MFU 约 0.5 已经是很强的实现}
课程给出的经验阈值是 MFU \(\ge 0.5\) 通常已经相当不错。这里的“不错”不是说剩余一半算力一定可追回，而是硬件峰值来自理想 shape、dtype 和 kernel；真实训练还要承受 launch、layout、非 matmul 算子与软件开销。若 MFU 明显低于这一量级，才应继续分解是算子形状、memory movement、通信还是 input pipeline 在拖慢。
\end{knowledgebox}

\subsection{本章小结}

Compute accounting 的步骤是：从 shape 算 FLOPs，用 benchmark 得实际时间，再与硬件峰值比较。不要只数 FLOPs，还要问这些 FLOPs 是否能高效喂给硬件。

